% ============================================================================
% De Modelos Sintéticos a Dados Reais — Validação Empírica com Datasets Públicos
% Beamer presentation – Portuguese (pt-BR)
% ============================================================================
\documentclass[aspectratio=169, 11pt]{beamer}

% ---------- Encoding & Language ----------
\usepackage[utf8]{inputenc}
\usepackage[T1]{fontenc}
\usepackage[brazil]{babel}

% ---------- Theme ----------
\usetheme{Madrid}
\usecolortheme{whale}
\setbeamertemplate{navigation symbols}{}
\setbeamertemplate{footline}{%
  \leavevmode%
  \hbox{%
    \begin{beamercolorbox}[wd=.333\paperwidth,ht=2.5ex,dp=1ex,center]{author in head/foot}%
      \usebeamerfont{author in head/foot}\insertshortauthor
    \end{beamercolorbox}%
    \begin{beamercolorbox}[wd=.333\paperwidth,ht=2.5ex,dp=1ex,center]{title in head/foot}%
      \usebeamerfont{title in head/foot}\insertshorttitle
    \end{beamercolorbox}%
    \begin{beamercolorbox}[wd=.333\paperwidth,ht=2.5ex,dp=1ex,right]{date in head/foot}%
      \usebeamerfont{date in head/foot}\insertshortdate{}\hspace*{2em}
      \insertframenumber{} / \inserttotalframenumber\hspace*{2ex}
    \end{beamercolorbox}}%
  \vskip0pt%
}

% ---------- Packages ----------
\usepackage{amsmath, amssymb}
\usepackage{booktabs}
\usepackage{graphicx}
\usepackage{tikz}
\usetikzlibrary{arrows.meta, positioning, shapes.geometric, calc}
\usepackage{multirow}
\usepackage{colortbl}
\usepackage{xcolor}
\usepackage{hyperref}

% ---------- Custom colors ----------
\definecolor{accent}{RGB}{37, 99, 235}
\definecolor{greenok}{RGB}{22, 163, 74}
\definecolor{redbad}{RGB}{220, 38, 38}
\definecolor{lightbg}{RGB}{248, 250, 252}
\definecolor{midblue}{RGB}{59, 130, 246}

% ---------- Meta ----------
\title[Modelos Sintéticos → Dados Reais]{%
  De modelos sintéticos a dados reais:\\[4pt]
  \large validação empírica em redes de sensores}
\author{Sensor Crossover \quad$\bullet$\quad Sunrise CPD}
\date{2026}
\institute{}

% ============================================================================
\begin{document}

% ============================================================================
% SLIDE 1 --- Título
% ============================================================================
\begin{frame}[plain]
\titlepage
\vspace{-1.5em}
\begin{center}
\small
\textcolor{gray}{%
  Modelos sintéticos guiam a teoria $\bullet$
  dados reais testam suas hipóteses}
\end{center}
\end{frame}

% ============================================================================
% SLIDE 2 --- Motivação original: artigo rejeitado
% ============================================================================
\begin{frame}{Motivação original --- o artigo rejeitado}

\begin{columns}[T]
\begin{column}{0.52\textwidth}

\textbf{Ponto de partida}
\begin{itemize}
  \item Linha teórica sobre aprendizado com \textbf{poucas amostras}
        e \textbf{poucas \textit{features}}.
  \item Modelos gaussianos sintéticos com classificadores lineares / SVM.
  \item Submissão a periódico: \textbf{rejeitada}~\textcolor{redbad}{$\times$}.
\end{itemize}

\vspace{0.4em}
\textbf{O que a rejeição indicou}
\begin{itemize}
  \item A crítica \emph{não} invalida o modelo.
  \item Aponta uma \textbf{fragilidade}: falta de validação em
        dados reais.
\end{itemize}

\end{column}
\begin{column}{0.44\textwidth}

\begin{block}{Mensagem}
\centering
O objetivo agora é transformar uma\\
\textbf{crítica dos revisores} em\\
uma \textbf{agenda experimental}.
\end{block}

\end{column}
\end{columns}

\end{frame}

% ============================================================================
% SLIDE 3 --- O que os revisores cobravam
% ============================================================================
\begin{frame}{O que os revisores cobravam}

\begin{center}
\small
\begin{tabular}{@{}p{5.0cm}p{7.0cm}@{}}
  \toprule
  \textbf{Crítica} & \textbf{Resposta do projeto} \\
  \midrule
  Falta de dados reais &
    Usar \textbf{Intel Lab} (Sensor Crossover) e
    \textbf{SensorScope} (Sunrise CPD). \\
  \addlinespace[3pt]
  Modelo muito idealizado &
    Comparar parâmetros sintéticos com parâmetros
    \textbf{empíricos} estimados nos datasets. \\
  \addlinespace[3pt]
  Motivação em sensores pouco explorada &
    Trabalhar explicitamente com \textbf{redes de sensores}
    e custos de coleta/transmissão. \\
  \addlinespace[3pt]
  Posicionamento na literatura &
    Conectar com detecção de mudança em redes
    multi-sensor com restrição de recursos. \\
  \bottomrule
\end{tabular}
\end{center}

\end{frame}

% ============================================================================
% SLIDE 4 --- Hipótese geral
% ============================================================================
\begin{frame}{Hipótese geral do projeto}

\begin{columns}[T]
\begin{column}{0.50\textwidth}

\textbf{Modelo sintético}
\begin{itemize}
  \item Controle de média, variância, covariância e correlação.
  \item Distribuição assumida (gaussiana).
  \item Reprodutibilidade exata.
\end{itemize}

\end{column}
\begin{column}{0.46\textwidth}

\textbf{Dado real}
\begin{itemize}
  \item Ruído e \textit{outliers}.
  \item Amostragem irregular e dados faltantes.
  \item Correlação empírica desconhecida.
  \item Distribuição não gaussiana.
\end{itemize}

\end{column}
\end{columns}

\vspace{0.8em}
\begin{block}{Pergunta central}
\centering
\itshape
Os fenômenos previstos por modelos idealizados continuam aparecendo em dados reais?
\end{block}

\end{frame}

% ============================================================================
% SLIDE 5 --- Mapa da agenda experimental
% ============================================================================
\begin{frame}{Mapa da agenda experimental}

\begin{center}
\begin{tikzpicture}[
    node distance=0.9cm and 0.9cm,
    eixo/.style={rectangle, draw=accent, fill=accent!8, rounded corners=4pt,
                 minimum height=1.2cm, minimum width=3.6cm, align=center,
                 font=\small},
  ]
  \node[eixo] (slacgs) {\textbf{SLACGS}\\modelo sintético\\(teoria)};
  \node[eixo, right=of slacgs] (sc) {\textbf{Sensor Crossover}\\Intel Lab\\(dados reais)};
  \node[eixo, right=of sc] (sun) {\textbf{Sunrise CPD}\\SensorScope\\(dados reais)};
\end{tikzpicture}
\end{center}

\vspace{0.6em}
\begin{itemize}
  \item \textbf{SLACGS:} estudo teórico do trade-off \textit{features}~$\times$~amostras.
  \item \textbf{Sensor Crossover:} validação empírica com Intel Lab.
  \item \textbf{Sunrise CPD:} adaptação de detecção de ponto de mudança em
        rede multi-sensor sob restrição de orçamento.
\end{itemize}

\vspace{0.4em}
\begin{center}
\itshape\small
Os dois projetos reais não \emph{substituem} o modelo sintético;\\
eles \textbf{testam e refinam} sua interpretação.
\end{center}

\end{frame}

% ============================================================================
% SLIDE 6 --- SLACGS: papel do modelo sintético
% ============================================================================
\begin{frame}{SLACGS --- papel do modelo sintético}

\begin{columns}[T]
\begin{column}{0.50\textwidth}

\textbf{O simulador permite controlar}
\begin{itemize}
  \item Médias das classes $\mu_0, \mu_1$.
  \item Variâncias $\sigma_i^2$.
  \item Covariância / correlação $\Sigma$.
  \item Número de \textit{features} $d$.
  \item Número de amostras $n$.
\end{itemize}

\end{column}
\begin{column}{0.46\textwidth}

\textbf{Fenômeno de interesse}
\begin{itemize}
  \item Mais \textit{features} podem \textbf{ajudar ou prejudicar}
        dependendo de $n$.
  \item Existe um ponto $n^{\star}$ em que usar
        mais \textit{features} passa a compensar.
\end{itemize}

\vspace{0.4em}
\begin{block}{Conceito}
\centering
\textit{Crossover:} ponto $n^{\star}$ a partir do qual
acrescentar uma \textit{feature} reduz o erro.
\end{block}

\end{column}
\end{columns}

\end{frame}

% ============================================================================
% SLIDE 7 --- Covariância, correlação e ponte com dados reais
% ============================================================================
\begin{frame}{Covariância, correlação e a ponte com dados reais}

\begin{columns}[T]
\begin{column}{0.50\textwidth}

\textbf{No modelo sintético}
\begin{itemize}
  \item A matriz $\Sigma$ é \emph{entrada} do modelo.
  \item Ela é \textbf{condicionada} à estrutura de classes.
\end{itemize}

\vspace{0.4em}
\textbf{Nos dados reais começamos observando}
\begin{itemize}
  \item Covariância empírica \textbf{global}.
  \item Correlação empírica global.
  \item Versões padronizadas.
\end{itemize}

\end{column}
\begin{column}{0.46\textwidth}

\begin{block}{Foi necessário distinguir}
\begin{enumerate}
  \item Covariância \textbf{teórica} do simulador.
  \item Covariância \textbf{empírica global} do dataset.
  \item Correlação entre sensores reais.
\end{enumerate}
\end{block}

\vspace{0.3em}
\small\itshape
A tentativa de validar com dados reais nos obrigou a entender melhor o que a covariância do modelo realmente representa.

\end{column}
\end{columns}

\end{frame}

% ============================================================================
% SLIDE 8 --- Sensor Crossover: ideia
% ============================================================================
\begin{frame}{Sensor Crossover --- ideia do experimento}

\begin{columns}[T]
\begin{column}{0.52\textwidth}

\textbf{Dataset:} Intel Berkeley Research Lab
\begin{itemize}
  \item 54 sensores de temperatura.
  \item ${\sim}\,2{,}3$ milhões de leituras em 38 dias.
\end{itemize}

\vspace{0.4em}
\textbf{Construção do problema}
\begin{itemize}
  \item Sensores de temperatura tratados como \textbf{features}.
  \item Um sensor de \textbf{referência} $R$ define as classes
        por separação na mediana.
  \item Outros sensores entram como \textit{features} adicionais.
\end{itemize}

\end{column}
\begin{column}{0.44\textwidth}

\begin{block}{Pergunta}
\centering
Mais sensores como \textit{features}\\
\textbf{sempre ajudam}?
\end{block}

\vspace{0.4em}
\small
Predição teórica: existe um ponto de cruzamento $n^{\star}$ que depende da correlação entre os sensores.

\end{column}
\end{columns}

\end{frame}

% ============================================================================
% SLIDE 9 --- Sensor Crossover: desenho experimental
% ============================================================================
\begin{frame}{Sensor Crossover --- desenho experimental}

\begin{columns}[T]
\begin{column}{0.50\textwidth}

\textbf{Sensores escolhidos}
\begin{itemize}
  \item $R = $ mote 22 (referência).
  \item $A = $ mote 25 (\textit{feature} primária).
  \item $B_{\text{alt}} = 31$, $B_{\text{med}} = 21$, $B_{\text{baixa}} = 48$
        (correlações decrescentes).
\end{itemize}

\vspace{0.4em}
\textbf{Classificador}
\begin{itemize}
  \item SVM linear, $C = 1{,}0$.
\end{itemize}

\end{column}
\begin{column}{0.46\textwidth}

\textbf{Procedimento}
\begin{itemize}
  \item Variar $n \in \{2, 4, \ldots, 1024\}$.
  \item 500 repetições Monte Carlo por $(n, d)$.
  \item Comparar erro com 1 e com 2 \textit{features}.
\end{itemize}

\vspace{0.4em}
\textbf{Métrica de crossover}
\[
  \Delta(n) = E_{d=1}(n) - E_{d=2}(n)
\]
$n^{\star}$: troca de sinal de $\Delta$.

\end{column}
\end{columns}

\end{frame}

% ============================================================================
% SLIDE 10 --- Sensor Crossover: resultados
% ============================================================================
\begin{frame}{Sensor Crossover --- resultados em dados reais}

\begin{center}
\small
\begin{tabular}{@{}lccc@{}}
  \toprule
  \textbf{Cenário} & $\boldsymbol{\rho}(B,R)$ & $\mathbf{n^{\star}}$ & \textbf{Tipo} \\
  \midrule
  Alta correlação  & 0{,}89 & 106     & Simples \\
  Média correlação & 0{,}66 & 4 e 77  & \textcolor{redbad}{\textbf{Duplo}} \\
  Baixa correlação & 0{,}38 & 14      & Simples \\
  \bottomrule
\end{tabular}
\end{center}

\vspace{0.6em}
\begin{exampleblock}{Mensagem}
\begin{itemize}
  \item O \textbf{crossover apareceu} em dados reais.
  \item A \textbf{correlação entre sensores} desloca $n^{\star}$.
  \item O dado real trouxe um caso de \textbf{cruzamento duplo},
        comportamento mais complexo do que o previsto inicialmente
        no simulador.
\end{itemize}
\end{exampleblock}

\end{frame}

% ============================================================================
% SLIDE 11 --- Sensor Crossover: aprendizado
% ============================================================================
\begin{frame}{Sensor Crossover --- aprendizado}

\begin{itemize}
  \item O experimento real reforça que ``mais \textit{features}''
        \textbf{não é automaticamente melhor}.
  \item A \textbf{qualidade} e a \textbf{relação estatística}
        entre sensores importam.
  \item Dados reais introduzem efeitos que \textbf{não estavam explícitos}
        no simulador (cauda das distribuições, regimes de operação,
        cruzamento duplo).
  \item Isso responde \emph{parcialmente} à crítica dos revisores:
        o fenômeno teórico sobreviveu à transição para dados reais.
\end{itemize}

\end{frame}

% ============================================================================
% SLIDE 12 --- Segunda frente: detecção de mudança
% ============================================================================
\begin{frame}{Segunda frente --- detecção de ponto de mudança}

\begin{columns}[T]
\begin{column}{0.52\textwidth}

\textbf{Mudança de problema}
\begin{itemize}
  \item De \emph{classificação} para \emph{detecção}.
  \item Tema: detecção de mudança em rede multi-sensor com
        \textbf{restrição de recursos}.
  \item Modelo teórico usa estatística \textbf{CUSUM} e
        política de amostragem sob orçamento.
\end{itemize}

\end{column}
\begin{column}{0.44\textwidth}

\textbf{Modelo}
\[
  X(t) \sim
  \begin{cases}
    \mathcal{N}(\mu_0, \sigma^2) & t < \tau \\
    \mathcal{N}(\mu_1, \sigma^2) & t \geq \tau
  \end{cases}
\]

\vspace{0.3em}
\small
$\tau$: instante da mudança.\\
$S(t)$: subconjunto de sensores escolhido dinamicamente.

\end{column}
\end{columns}

\end{frame}

% ============================================================================
% SLIDE 13 --- Sunrise CPD: adaptação para dados reais
% ============================================================================
\begin{frame}{Sunrise CPD --- adaptação para dados reais}

\begin{columns}[T]
\begin{column}{0.54\textwidth}

\textbf{Dataset:} SensorScope --- Grand-St-Bernard (EPFL)
\begin{itemize}
  \item Rede de sensores ambientais sem fio em alta montanha.
  \item Variável usada: \textbf{radiação solar}.
  \item Período: setembro--outubro de 2007 (43 dias válidos).
\end{itemize}

\vspace{0.4em}
\textbf{Ground truth externo}
\begin{itemize}
  \item Nascer do sol \textbf{astronômico}.
  \item Evento diário com horário conhecido externamente.
\end{itemize}

\end{column}
\begin{column}{0.42\textwidth}

\begin{block}{Ideia}
Detectar o nascer do sol como ponto de mudança \textbf{real}, com
ground truth independente do detector.
\end{block}

\end{column}
\end{columns}

\end{frame}

% ============================================================================
% SLIDE 14 --- Sunrise CPD: pipeline implementado
% ============================================================================
\begin{frame}{Sunrise CPD --- pipeline implementado}

\begin{columns}[T]
\begin{column}{0.50\textwidth}

\textbf{Etapas}
\begin{enumerate}
  \item Aquisição e pré-processamento do SensorScope.
  \item Relatório HTML do dataset.
  \item Cálculo do \emph{ground truth} astronômico.
  \item Ranking global de sensores por divergência gaussiana $D_i$.
  \item Detector \textbf{CUSUM} sobre razão de verossimilhança.
  \item Política dinâmica de seleção de sensores.
\end{enumerate}

\end{column}
\begin{column}{0.46\textwidth}

\textbf{Saídas}
\begin{itemize}
  \item Relatório HTML por cenário:
    \begin{itemize}
      \item \textit{low budget};
      \item \textit{medium budget};
      \item \textit{high budget}.
    \end{itemize}
  \item Relatório de \textbf{comparação} entre os cenários.
  \item JSONs com métricas agregadas por dia.
\end{itemize}

\end{column}
\end{columns}

\end{frame}

% ============================================================================
% SLIDE 15 --- Sunrise CPD: budget e seleção dinâmica
% ============================================================================
\begin{frame}{Sunrise CPD --- \textit{budget} e seleção dinâmica}

\begin{columns}[T]
\begin{column}{0.52\textwidth}

\textbf{Duas restrições} (do artigo)
\[
  \sum_{i \in S(t)} C_i \;\leq\; B
  \qquad
  \sum_{i \in S(t)} T_i \;\leq\; C
\]

\begin{itemize}
  \item $C_i$: custo de \textbf{coleta} (sensoriamento).
  \item $T_i$: custo de \textbf{transmissão}.
  \item $S(t)$: subconjunto escolhido a cada instante.
\end{itemize}

\vspace{0.4em}
\textbf{Custos sintéticos homogêneos:} $C_i = T_i = 1$.\\
O primeiro sensor escolhido é \textbf{referência local}
(custo efetivo de transmissão $0$); os demais são
\textbf{cooperativos}.

\end{column}
\begin{column}{0.42\textwidth}

\textbf{Regimes}

\vspace{0.3em}
\small
\begin{tabular}{@{}lcc@{}}
  \toprule
  \textbf{Regime} & $B$ & $C$ \\
  \midrule
  \textit{Low}    & 1 & 0 \\
  \textit{Medium} & 3 & 2 \\
  \textit{High}   & 9 & 8 \\
  \bottomrule
\end{tabular}

\vspace{0.7em}
\itshape
Percebemos que ``\textit{budget}'' não podia ser apenas um conjunto fixo de sensores.

\end{column}
\end{columns}

\end{frame}

% ============================================================================
% SLIDE 16 --- Sunrise CPD: resultados atuais
% ============================================================================
\begin{frame}{Sunrise CPD --- resultados atuais}

\begin{center}
\small
\begin{tabular}{@{}lccccc@{}}
  \toprule
  \textbf{Regime} & $|S(t)|$ médio &
  \textbf{No prazo} & \textbf{Atrasadas} & \textbf{Não detectadas} &
  \textbf{Atraso médio$^{+}$} \\
  \midrule
  \textit{Low}    & 0{,}97 & 0/43 & 39/43 & 4/43 & ${\sim}\,66$ min \\
  \textit{Medium} & 2{,}90 & 0/43 & 41/43 & 2/43 & ${\sim}\,69$ min \\
  \textit{High}   & 6{,}62 & 0/43 & 41/43 & 2/43 & ${\sim}\,78$ min \\
  \bottomrule
\end{tabular}
\end{center}

\vspace{0.3em}
\small
Tolerância: $\pm 15$ min em torno do nascer do sol astronômico.\\
``Atraso médio$^{+}$'' é a média de $\max(\tau - \nu, 0)$ nas detecções (estilo ADD).

\vspace{0.5em}
\begin{block}{Mensagem}
\begin{itemize}
  \item A \textbf{cooperação operacional} acontece: mais orçamento $\Rightarrow$ mais sensores ativos.
  \item Mas a detecção ficou \textbf{sistematicamente atrasada}.
  \item \textbf{Nenhum} regime ficou dentro da tolerância.
  \item Aumentar o orçamento \textbf{não melhorou a pontualidade}.
\end{itemize}
\end{block}

\end{frame}

% ============================================================================
% SLIDE 17 --- Sunrise CPD: o que deu errado e por quê
% ============================================================================
\begin{frame}{Sunrise CPD --- o que deu errado e por quê}

\begin{itemize}
  \item O modelo com parâmetros gaussianos \textbf{globais} parece
        inadequado para o nascer do sol real.
  \item O nascer do sol \textbf{não é uma mudança abrupta};
        é uma transição gradual.
  \item $\mu_1$ global pode representar a radiação \textbf{bem depois}
        do nascer do sol.
  \item O CUSUM acaba detectando quando o aumento já está \textbf{consolidado}.
  \item A média dos sensores em $S(t)$ pode \textbf{diluir} sinais fortes
        e atrasar a estatística.
\end{itemize}

\vspace{0.6em}
\begin{exampleblock}{Mensagem}
O problema deixou de ser apenas o \textit{budget}; passou a ser também
a \textbf{parametrização estatística} do evento real.
\end{exampleblock}

\end{frame}

% ============================================================================
% SLIDE 18 --- Discussão metodológica: global, diário e futuro
% ============================================================================
\begin{frame}{Discussão metodológica --- global, diário e futuro}

\begin{columns}[T]
\begin{column}{0.50\textwidth}

\textbf{Parâmetros globais}
\begin{itemize}
  \item $D_i$ global: útil como \emph{ranking}, mas usa
        informação agregada do dataset.
  \item $\mu_0$ global: problemático, porque cada dia tem
        \textbf{baseline e horário diferentes}.
\end{itemize}

\vspace{0.4em}
\textbf{Parâmetros diários}
\begin{itemize}
  \item $\mu_0$ diário: defensável (usa só pré-\textit{sunrise}).
  \item $D_i$ diário: \emph{quase oracle} (exige pós-\textit{sunrise}).
\end{itemize}

\end{column}
\begin{column}{0.46\textwidth}

\textbf{Caminhos futuros}
\begin{itemize}
  \item Separar \textbf{treino e teste} de forma temporal.
  \item Estimar parâmetros com dados \textbf{históricos}.
  \item Combinar $\mu_0$ diário com $D_i$ global.
  \item \textbf{Clusterizar dias} por padrões de visibilidade
        e radiação; estimar parâmetros por grupo.
\end{itemize}

\end{column}
\end{columns}

\end{frame}

% ============================================================================
% SLIDE 19 --- Síntese dos aprendizados
% ============================================================================
\begin{frame}{Síntese dos aprendizados}

\begin{center}
\small
\begin{tabular}{@{}lp{3.6cm}p{5.2cm}@{}}
  \toprule
  \textbf{Projeto} & \textbf{Resultado} & \textbf{Aprendizado} \\
  \midrule
  SLACGS &
    Modelo controlado &
    Esclarece o papel de covariância e correlação. \\
  \addlinespace[2pt]
  Sensor Crossover &
    Crossover observado em dados reais &
    O fenômeno sintético aparece, com complexidade adicional. \\
  \addlinespace[2pt]
  Sunrise CPD &
    \textit{Budget} dinâmico implementado, mas detecção atrasada &
    O dado real revela um \emph{mismatch} entre modelo e fenômeno. \\
  \bottomrule
\end{tabular}
\end{center}

\vspace{0.8em}
\begin{block}{Mensagem}
\centering\itshape
O dado real não é apenas validação; ele é um \textbf{instrumento
para descobrir hipóteses escondidas}.
\end{block}

\end{frame}

% ============================================================================
% SLIDE 20 --- Etapas futuras (inclui adversarial)
% ============================================================================
\begin{frame}{Etapas futuras}

\begin{columns}[T]
\begin{column}{0.50\textwidth}

\textbf{Sensor Crossover}
\begin{itemize}
  \item Consolidar como \textbf{evidência empírica} do primeiro artigo.
  \item Estender para mais combinações de sensores e dimensões.
\end{itemize}

\vspace{0.4em}
\textbf{Sunrise CPD}
\begin{itemize}
  \item Testar baselines \textbf{diários}.
  \item Calibrar limiar e \textit{drift} do CUSUM.
  \item Avaliar agregação ponderada e custos heterogêneos.
  \item Separar treino e teste no tempo.
\end{itemize}

\end{column}
\begin{column}{0.46\textwidth}

\textbf{Expansão para cenários adversariais}
\begin{itemize}
  \item Sensores comprometidos.
  \item Mudanças induzidas (ataques sutis).
  \item Cooperação sob \textbf{restrição de comunicação}.
  \item Conexão com detecção \textbf{robusta} em redes de sensores
        e com a parte covert/adversarial do artigo original.
\end{itemize}

\end{column}
\end{columns}

\end{frame}

% ============================================================================
% SLIDE 21 --- Fechamento
% ============================================================================
\begin{frame}{Fechamento}

\begin{itemize}
  \item Já há \textbf{progresso concreto} em duas frentes empíricas.
  \item O \textbf{primeiro experimento} (Sensor Crossover) produziu
        evidência positiva: o crossover aparece em dados reais.
  \item O \textbf{segundo experimento} (Sunrise CPD) revelou desafios
        metodológicos importantes na transição teoria $\to$ dado real.
  \item O principal ganho foi transformar uma crítica
        (\textit{``faltam dados reais''}) em uma agenda experimental
        com resultados verificáveis.
\end{itemize}

\vspace{1.0em}
\begin{center}
\begin{minipage}{0.88\textwidth}
\centering\itshape
\textcolor{accent}{\large
``Experimentar com dados reais não apenas testa a teoria;\\
muitas vezes ajuda a entender melhor a própria teoria.''}
\end{minipage}
\end{center}

\end{frame}

\end{document}
